\chapter{Dynamic programming}
\label{ch:dp}

When the dynamics $p$ of the MDP are known, the Bellman equations of
\cref{ch:mdp} can be solved by iterating the Bellman operators: this is
\emph{dynamic programming} (DP). DP is rarely applicable as such (the
dynamics are often unknown, and the number of states is usually huge, as in
Pacman), but every RL algorithm of this course is an approximation of one of
its two schemes, policy iteration and value iteration. This chapter
presents them with complete convergence proofs and solves exactly a
reduced version of Pacman.

\section{Policy evaluation}

\emph{Policy evaluation} computes $v_\pi$ for a given policy $\pi$. By
\cref{thm:bellman-exp} it is the unique fixed point of the contraction
$\T^\pi$, so the iteration $v_{k+1} = \T^\pi v_k$ converges to it from any
starting point.

\begin{algobox}{Iterative policy evaluation}
  \begin{algorithmic}[1]
    \Require policy $\pi$, tolerance $\varepsilon > 0$
    \State $v(s) \gets 0$ for all $s$
    \Repeat
      \State $\Delta \gets 0$
      \For{each state $s$}
        \State $u \gets \sum_a \pi(a \mid s) \sum_{s', r} p(s', r \mid s, a)
          \left[ r + \gamma\, v(s') \right]$
        \State $\Delta \gets \max(\Delta, \abs{u - v(s)})$;\quad $v(s) \gets u$
      \EndFor
    \Until{$\Delta < \varepsilon$}
    \State \Return $v$
  \end{algorithmic}
\end{algobox}

The algorithm above updates $v$ \emph{in place}: the new value of a state is
used immediately for the following states of the same sweep. This variant
(Gauss--Seidel) also converges, usually faster; we prove the convergence of
the synchronous variant, which is a direct consequence of the contraction.

\begin{proposition}{Convergence of policy evaluation}{pe-conv}
  Let $v_{k+1} = \T^\pi v_k$. Then $\norm{v_k - v_\pi}_\infty \leq
  \gamma^k \norm{v_0 - v_\pi}_\infty$, and if $\norm{v_{k+1} -
  v_k}_\infty \leq \varepsilon$ then
  \begin{equation}
    \norm{v_{k+1} - v_\pi}_\infty \leq \frac{\gamma\, \varepsilon}{1 -
    \gamma}.
  \end{equation}
\end{proposition}

\begin{proof}
  The first inequality follows by induction from $\norm{v_{k+1} -
  v_\pi}_\infty = \norm{\T^\pi v_k - \T^\pi v_\pi}_\infty \leq \gamma
  \norm{v_k - v_\pi}_\infty$. For the second,
  \begin{align*}
    \norm{v_{k+1} - v_\pi}_\infty
    &= \norm{\T^\pi v_k - \T^\pi v_\pi}_\infty
      \leq \gamma \norm{v_k - v_\pi}_\infty \\
    &\leq \gamma \left( \norm{v_k - v_{k+1}}_\infty
      + \norm{v_{k+1} - v_\pi}_\infty \right)
      \leq \gamma\, \varepsilon + \gamma \norm{v_{k+1} - v_\pi}_\infty,
  \end{align*}
  and we solve for $\norm{v_{k+1} - v_\pi}_\infty$.
\end{proof}

\begin{intuition}
  A small change between two sweeps does not mean that the values are
  almost right: each sweep only reduces the error by the factor $\gamma$,
  so the remaining error can be $\gamma / (1 - \gamma)$ times the last
  change, $99$ times with $\gamma = 0.99$. Slowly changing values may still
  be far from the answer, and the stopping tolerance must be chosen
  accordingly.
\end{intuition}

\section{Policy improvement}

Knowing $v_\pi$, can we find a better policy? The natural candidate is the
policy that is greedy with respect to $q_\pi$: in each state, play the
action that is best if $\pi$ is followed afterwards.

\begin{theorem}{Policy improvement}{pi-improvement}
  Let $\pi$ and $\pi'$ be two policies such that, for every state $s$,
  \begin{equation}
    \label{eq:improvement-cond}
    \sum_a \pi'(a \mid s)\, q_\pi(s, a) \geq v_\pi(s).
  \end{equation}
  Then $v_{\pi'} \geq v_\pi$. If moreover the inequality
  \eqref{eq:improvement-cond} is strict in some state $s$, then
  $v_{\pi'}(s) > v_\pi(s)$.
\end{theorem}

\begin{proof}
  The left-hand side of \eqref{eq:improvement-cond} is $(\T^{\pi'}
  v_\pi)(s)$, by \eqref{eq:q-from-v}. The hypothesis is therefore
  $\T^{\pi'} v_\pi \geq v_\pi$. By monotonicity of $\T^{\pi'}$
  (\cref{lem:contraction}) and induction, $(\T^{\pi'})^n v_\pi \geq
  (\T^{\pi'})^{n-1} v_\pi \geq \dots \geq v_\pi$ for every $n$, and
  $(\T^{\pi'})^n v_\pi \to v_{\pi'}$ as $n \to \infty$. Hence
  $v_{\pi'} \geq \T^{\pi'} v_\pi \geq v_\pi$, the first inequality
  being the limit of $(\T^{\pi'})^n v_\pi \geq \T^{\pi'} v_\pi$. If
  $(\T^{\pi'} v_\pi)(s) > v_\pi(s)$, then $v_{\pi'}(s) \geq
  (\T^{\pi'} v_\pi)(s) > v_\pi(s)$.
\end{proof}

\begin{intuition}
  The hypothesis compares two plans that differ only in the first move:
  play $\pi'$ once, then return to $\pi$, versus play $\pi$ throughout. If
  deviating for one move never hurts, then deviating at every move does not
  hurt either: each deviation starts from a state where it is at least as
  good, and the gains add up along the trajectory.
\end{intuition}

\begin{corollary}{Greedy improvement}{greedy-improvement}
  Let $\pi'$ be greedy with respect to $q_\pi$: $\pi'(s) \in \argmax_a
  q_\pi(s, a)$ for every $s$. Then $v_{\pi'} \geq v_\pi$. If $v_{\pi'} =
  v_\pi$, then $\pi$ and $\pi'$ are both optimal.
\end{corollary}

\begin{proof}
  $q_\pi(s, \pi'(s)) = \max_a q_\pi(s, a) \geq \sum_a \pi(a \mid s)\,
  q_\pi(s, a) = v_\pi(s)$, so \cref{thm:pi-improvement} applies. If
  $v_{\pi'} = v_\pi$, then by \eqref{eq:q-from-v} $q_{\pi'} = q_\pi$, and
  \begin{equation*}
    v_{\pi'}(s) = q_{\pi'}(s, \pi'(s)) = q_\pi(s, \pi'(s))
    = \max_a q_\pi(s, a)
    = \max_a \left[ r(s, a) + \gamma \sum_{s'} p(s' \mid s, a)\,
      v_{\pi'}(s') \right]:
  \end{equation*}
  $v_{\pi'}$ is a fixed point of $\T^\star$, hence equals $v_\star$
  (\cref{thm:bellman-opt}).
\end{proof}

\begin{intuition}
  To improve a policy it suffices to ask, in every state, ``which action
  would be best if I followed my current policy afterwards?'', and to play
  it. If the answer is always the action already played, no one-move
  deviation helps, and then no deviation of any kind helps: the policy is
  optimal. Optimality can be checked one move at a time.
\end{intuition}

\section{Policy iteration}

Alternating evaluation and greedy improvement gives \emph{policy
iteration}.

\begin{algobox}{Policy iteration}
  \begin{algorithmic}[1]
    \State choose any deterministic policy $\pi_0$
    \For{$k = 0, 1, 2, \dots$}
      \State \textbf{evaluation:} compute $v_{\pi_k}$
        (\cref{prop:matrix} or iterative policy evaluation)
      \State \textbf{improvement:} $\pi_{k+1}(s) \gets \argmax_a \left[
        r(s, a) + \gamma \sum_{s'} p(s' \mid s, a)\, v_{\pi_k}(s') \right]$,
        keeping $\pi_k(s)$ in case of a tie
      \If{$\pi_{k+1} = \pi_k$} \Return $\pi_k$, $v_{\pi_k}$ \EndIf
    \EndFor
  \end{algorithmic}
\end{algobox}

\begin{theorem}{Convergence of policy iteration}{pi-conv}
  Policy iteration stops after a finite number of iterations, at most the
  number $\abs{\Actions}^{\abs{\States}}$ of deterministic policies, and
  returns an optimal policy.
\end{theorem}

\begin{proof}
  If $\pi_{k+1} \neq \pi_k$, then because ties keep the previous action,
  $\pi_{k+1}(s)$ strictly improves on $\pi_k(s)$ in some state:
  $q_{\pi_k}(s, \pi_{k+1}(s)) > q_{\pi_k}(s, \pi_k(s)) = v_{\pi_k}(s)$. By
  \cref{thm:pi-improvement}, $v_{\pi_{k+1}} \geq v_{\pi_k}$ with strict
  inequality in $s$. The value functions $v_{\pi_0}, v_{\pi_1}, \dots$
  therefore strictly increase (for the partial order of functions), so no
  policy appears twice; since there are finitely many deterministic
  policies, the algorithm stops. When it stops, $\pi_k$ is greedy with
  respect to its own $q_{\pi_k}$, so $\T^\star v_{\pi_k} = \T^{\pi_k}
  v_{\pi_k} = v_{\pi_k}$: $v_{\pi_k} = v_\star$ and $\pi_k$ is optimal.
\end{proof}

The bound $\abs{\Actions}^{\abs{\States}}$ is very pessimistic: in practice
policy iteration converges in a handful of iterations, as the Pacman example
below shows.

\section{Value iteration}

Policy iteration evaluates each policy completely before improving it.
\emph{Value iteration} goes to the other extreme: it applies the Bellman
optimality operator directly, $v_{k+1} = \T^\star v_k$, which amounts to
one sweep of evaluation interleaved with an improvement in every state.

\begin{algobox}{Value iteration}
  \begin{algorithmic}[1]
    \Require tolerance $\varepsilon > 0$
    \State $v(s) \gets 0$ for all $s$
    \Repeat
      \State $v_{\text{old}} \gets v$
      \For{each state $s$}
        \State $v(s) \gets \max_a \left[ r(s, a) + \gamma \sum_{s'}
          p(s' \mid s, a)\, v_{\text{old}}(s') \right]$
      \EndFor
    \Until{$\norm{v - v_{\text{old}}}_\infty < \varepsilon$}
    \State \Return $v$ and a greedy policy with respect to $v$
  \end{algorithmic}
\end{algobox}

\begin{theorem}{Convergence of value iteration}{vi-conv}
  Let $v_{k+1} = \T^\star v_k$.
  \begin{enumerate}[nosep, label=(\roman*)]
    \item $\norm{v_k - v_\star}_\infty \leq \gamma^k \norm{v_0 -
      v_\star}_\infty$.
    \item If $\norm{v_{k+1} - v_k}_\infty \leq \varepsilon$, then
      $\norm{v_{k+1} - v_\star}_\infty \leq \gamma \varepsilon / (1 -
      \gamma)$.
    \item If $\norm{v - v_\star}_\infty \leq \delta$ and $\pi$ is greedy
      with respect to $v$, then $\norm{v_\pi - v_\star}_\infty \leq 2
      \gamma \delta / (1 - \gamma)$.
  \end{enumerate}
\end{theorem}

\begin{proof}
  (i) and (ii) are proved exactly as \cref{prop:pe-conv}, with $\T^\star$
  and its fixed point $v_\star$ in place of $\T^\pi$ and $v_\pi$.

  (iii) Since $\pi$ is greedy with respect to $v$, $\T^\pi v = \T^\star v$.
  Then
  \begin{align*}
    \norm{v_\pi - v_\star}_\infty
    &\leq \norm{\T^\pi v_\pi - \T^\pi v}_\infty
      + \norm{\T^\star v - \T^\star v_\star}_\infty
    \leq \gamma \norm{v_\pi - v}_\infty + \gamma \norm{v - v_\star}_\infty
    \\
    &\leq \gamma \left( \norm{v_\pi - v_\star}_\infty
      + \norm{v_\star - v}_\infty \right) + \gamma\, \delta
    \leq \gamma \norm{v_\pi - v_\star}_\infty + 2 \gamma\, \delta,
  \end{align*}
  using $v_\pi = \T^\pi v_\pi$ and $v_\star = \T^\star v_\star$ in the first
  line. Solving for $\norm{v_\pi - v_\star}_\infty$ gives the result.
\end{proof}

\begin{intuition}
  Statement (iii) is what matters in practice: approximate values give
  good decisions. An error $\delta$ in the values can make the agent choose
  an action that is worse by about $2 \delta$ (one action overestimated by
  $\delta$, the best one underestimated by $\delta$), and such losses can
  occur at every step, hence the factor of the horizon $1 / (1 - \gamma)$.
  To act well, the values need not be exact; they need to rank the actions
  correctly.
\end{intuition}

\paragraph{Asynchronous DP.} Nothing forces the backups to sweep all the
states in order. \emph{Asynchronous} value iteration updates the states in
any order, one at a time, using the latest values (as the in-place policy
evaluation above): it still converges to $v_\star$ as long as every state is
updated infinitely often (Bertsekas and Tsitsiklis, 1989). This freedom
allows focusing the computation on the states that matter, for example the
states visited by the current policy (\emph{real-time dynamic
programming}), or the states whose value changed most recently
(\emph{prioritised sweeping}, \cref{ch:model-based}).

\section{Linear programming}

The optimal value is also the solution of a linear program, which gives
another exact method and the basis of several modern algorithms.

\begin{theorem}{LP formulation}{lp}
  Let $\mu$ be any distribution with $\mu(s) > 0$ for all $s$. Then
  $v_\star$ is the unique solution of
  \begin{equation}
    \label{eq:lp}
    \min_{v \in \R^{\States}} \sum_s \mu(s)\, v(s) \quad
    \text{subject to} \quad v(s) \geq r(s, a) + \gamma \sum_{s'} p(s' \mid
    s, a)\, v(s') \quad \text{for all } s, a.
  \end{equation}
\end{theorem}

\begin{proof}
  The constraints say exactly $v \geq \T^\star v$. Then by monotonicity of
  $\T^\star$ (\cref{lem:contraction}), $v \geq \T^\star v \geq
  (\T^\star)^2 v \geq \dots$, and the iterates converge to $v_\star$
  (\cref{thm:bellman-opt}): every feasible $v$ satisfies $v \geq v_\star$.
  Conversely $v_\star$ is feasible ($v_\star = \T^\star v_\star$). Since
  $\mu > 0$, $\sum_s \mu(s) v(s) \geq \sum_s \mu(s) v_\star(s)$ with
  equality only if $v = v_\star$.
\end{proof}

\begin{intuition}
  A feasible $v$ is an \emph{optimistic} value: in every state it is at
  least what any action gives if the future is valued by $v$. Such a
  function overestimates $v_\star$ everywhere, and the linear program looks
  for the lowest one: the tightest ``roof'' over the one-step values of all
  the actions, which touches them exactly at $v_\star$.
\end{intuition}

The program has $\abs{\States}$ variables and $\abs{\States}
\abs{\Actions}$ constraints. Its dual has one variable $x(s, a) \geq 0$ per
state-action pair, which can be interpreted as the discounted frequency of
playing $a$ in $s$ (an \emph{occupancy measure}); optimal policies are read
from the optimal dual solution as $\pi(a \mid s) \propto x(s, a)$. This
dual view underlies several offline RL methods (\cref{ch:imitation}).

\begin{remark}{Generalised policy iteration}{gpi}
  Policy iteration and value iteration are two extremes of a common scheme,
  \emph{generalised policy iteration}: a value function is pushed towards
  the value of the current policy (evaluation), while the policy is pushed
  towards greediness with respect to the value function (improvement). The
  two processes compete (each changes the target of the other) and
  cooperate (their only joint fixed point is optimal). All the algorithms of
  the next chapters follow this pattern: Q-learning and DQN are
  approximations of value iteration, actor-critic methods and PPO of
  policy iteration.
\end{remark}

\section{Pacman solved exactly}
\label{sec:dp-pacman}

The full game is far too large for DP (\cref{sec:curse}), but a reduced
version can be solved exactly, which shows what an optimal policy does.

\begin{pacman}{Reaching the power cookie}{dp-escape}
  On the maze \code{pacman\_map3} of game \code{v1}, Pacman and \emph{one}
  ghost (aggressiveness $0.74$, chasing along shortest paths), no cookies.
  Pacman must reach a power cookie in the bottom right corner of the maze:
  reaching it gives $+10$ and ends the episode, being eaten gives $-35$,
  and each move costs $0.15$, with $\gamma = 0.99$. Pacman moves first,
  then the ghost. The state is the pair (Pacman's cell, ghost's cell):
  $77 \times 77 = 5929$ states. The script \code{Course/code/dp\_escape.py}
  builds the transition probabilities and solves the MDP.
  \begin{itemize}[nosep]
    \item Value iteration reaches $\norm{v_{k+1} - v_k}_\infty < 10^{-8}$
      after $57$ sweeps; policy iteration stops after $11$ iterations, far
      below the bound $5^{5929}$ of \cref{thm:pi-conv}. Both give the same
      policy, and value functions that differ by $10^{-8}$.
    \item The optimal policy reaches the power cookie \emph{from every
      starting state} with probability $1$ (to $10^{-12}$). The policy that
      follows a shortest path to the cookie, ignoring the ghost, succeeds
      from only $67.9\%$ of the states on average.
    \item From Pacman at $(5, 7)$ with the ghost at $(7, 12)$, between him
      and the cookie, the shortest path succeeds with probability $0.004$.
      The optimal policy succeeds with probability $1$, in $11.4$ moves on
      average instead of the $10$ moves of the shortest path: it makes a
      detour around the ghost. Its value is $v_\star = 7.39$.
  \end{itemize}
  The fact that the optimal success probability is $1$ everywhere is a
  property of the maze: it has no dead ends, and a single ghost that moves
  at Pacman's speed, after him, can never corner him. This is why the game
  needs (at least) two ghosts.
\end{pacman}

The heart of the script is the Bellman backup and the value iteration
loop; \code{T[k][a]} lists the triples (probability, reward, next state) of
the transition from state \code{k} with action \code{a}, the next state
being $-1$ for a terminal state.
\begin{lstlisting}[caption={Value iteration (\code{Course/code/dp\_escape.py},
  excerpt)}]
def backup(v, k, a):
    # r(s,a) + gamma * sum_s' p(s'|s,a) v(s'), terminal states worth 0
    return sum(pr * (r + (GAMMA * v[s2] if s2 >= 0 else 0.0))
               for pr, r, s2 in T[k][a])

def value_iteration(tol=1e-8):
    v = np.zeros(nS); it = 0
    while True:
        it += 1
        new = np.array([max(backup(v, k, a) for a in range(nA))
                        for k in range(nS)])      # v <- T* v
        delta = np.max(np.abs(new - v)); v = new
        if delta < tol:
            return v, it
\end{lstlisting}

\section{The curse of dimensionality}
\label{sec:curse}

DP needs two things that are missing in practice.
\begin{itemize}
  \item \textbf{The model.} Each backup uses $p(s', r \mid s, a)$. For
    Pacman, these probabilities could be derived from the code of the
    ghosts, but in general the environment is a black box that can only be
    sampled.
  \item \textbf{A table of all states.} Each sweep visits every state.
\end{itemize}

\begin{pacman}{How many states?}{n-states}
  In game \code{v1}, a state specifies Pacman's cell ($77$ choices), the
  cell of each ghost or the fact that it was eaten ($78$ choices each),
  which of the $77$ cells still hold a cookie (at most $2^{77}$ choices),
  and the power credit ($13$ values). This gives at most
  \begin{equation*}
    77 \times 78^2 \times 2^{77} \times 13 \approx 9 \times 10^{29}
  \end{equation*}
  states. Most of these configurations cannot occur, but even a tiny
  fraction of them is far beyond what a table can store: at one byte per
  state, $10^{20}$ states would already require $100$ exabytes.
\end{pacman}

The two reduced games of this chapter illustrate the jump: one ghost and no
cookies give $5929$ states, solved in seconds; adding the cookies
multiplies the number of states by up to $2^{77}$. The way out, developed in
the rest of the course, is twofold: \emph{learn from sampled experience}
instead of the model (\cref{ch:td}), and \emph{generalise} between states
with a parametrised function, a neural network reading the image of the
board (\cref{ch:dqn}).
